\subproblem{Үйкеліс арқылы дәнекерлеу}{3.5}

\textit{Үйкеліс арқылы дәнекерлеудің} қарапайым моделін қарастырайық. Радиусы $R=\qty1{\cm}$, осьтері сәйкес келетін екі бірдей жеткілікті ұзын болат цилиндрлер берілген. Цилиндрлердің бірі қандай да бір $\omega$ бұрыштық жылдамдығымен айналдырылады және ол бекітілген цилиндрдің ұшына $p=\qty{20}{\MPa}$ тұрақты қысыммен қысылады. Шекті $\omega$-ның қандай мәнінде болат стационарлы режимде (цилиндр жеткілікті ұзақ уақыт айналғаннан кейін) балқи бастайтынын анықтаңыз. Цилиндр радиусы бойынша температура тез теңеседі деп есептеңіз. Қоршаған орта температурасы $T_0=\qty{20}{\celsius}$. Тот баспайтын болат үшін: болаттың болатқа үйкеліс коэффициенті $\mu=0.4$, жылуөткізгіштік коэффициенті $\kappa=\qty{20}{\WmK}$, мәжбүрлі конвекция кезіндегі ауаға жылу беру коэффициенті $\alpha=\qty{200}{\W\per\metre\squared\per\kelvin}$, ал балқу нүктесі $T^*=\qty{1520}{\celsius}$. Болаттың параметрлерін оның температурасына тәуелсіз деп санаңыз; цилиндрлердің жылу сыйымдылығын ескермеңіз.

\textbf{Математикалық кеңес:} Қандай да бір нақты $\lambda$ заттық тұрақтысы үшін $x''(t) = \lambda^2 x(t)$ түріндегі дифференциалдық теңдеудің шешімі келесідей болады:
$$
x(t) = Ae^{\lambda t} + Be^{-\lambda t}
$$
мұндағы $A$ және $B$ --- интегралдау тұрақтылары, ал $e=2.718...$ --- Эйлер саны.
